\section{Search-Guided Recursive Decimated-BP Decoder}
\label{sec:search_guided_recursive_bp}

\subsection{目的と基本的な考え方}

本節では、Belief Propagation (BP) と correction search を組み合わせた
探索型復号器を定式化する。
本手法の目的は、単純に BP の不確かな variable node を順番に固定することではなく、
以下の二つの処理を同一の探索過程の中で同時に行うことである。

\begin{enumerate}
    \item 低コストの valid correction に直接到達する可能性の高い探索を行う。
    \item valid correction が探索予算内に得られない場合でも、
    次に実行する decimated BP に適した hard-fixation pattern を取得する。
\end{enumerate}

したがって、本手法では search node に二種類の役割を持たせる。
第一は correction search としての価値であり、
これを solve score $F_{\mathrm{solve}}$ により評価する。
第二は、次の BP instance を生成するための decimation pattern としての価値であり、
これを guide score $F_{\mathrm{guide}}$ により評価する。

この二つを分離することが本手法の基本的な設計方針である。
Tesseract~\cite{tesseract} は前者に相当する correction search を行う一方、
Beam Search Decoder~\cite{beamsearch} は BP posterior を用いて
decimation path を選択する。
本手法はこれらを直接混合するのではなく、
一つの BP state から多数の search node を低コストに評価し、
その中の少数のみについて実際に decimated BP を実行する階層構造を採用する。


\subsection{問題設定と記法}

binary decoding matrix を
\begin{equation}
    H \in \mathbb{F}_2^{M\times N}
\end{equation}
とし、観測 syndrome を
\begin{equation}
    \bm{s}\in\mathbb{F}_2^M
\end{equation}
とする。
variable node の集合を
\begin{equation}
    \mathcal{V}=\{1,\ldots,N\},
\end{equation}
check node の集合を
\begin{equation}
    \mathcal{C}=\{1,\ldots,M\}
\end{equation}
とする。

variable $j$ の physical error probability を $p_j$ とし、
対応する physical weight を
\begin{equation}
    w_j
    =
    \log\frac{1-p_j}{p_j}
\end{equation}
と定義する。

check $a$ に接続する variable 集合を
\begin{equation}
    \mathcal{N}(a)
    =
    \{j:H_{aj}=1\}
\end{equation}
とする。

本手法では複数の BP instance を保持する。
BP instance を $B$ と書き、その instance において既に hard-fix された
variable assignment の集合を
\begin{equation}
    D_B
    =
    \{(j,b_j):j\in\mathcal{F}_B,\;b_j\in\{0,1\}\}
\end{equation}
とする。
ここで $\mathcal{F}_B$ は fixed variable の集合である。
free variable 集合を
\begin{equation}
    \mathcal{U}_B
    =
    \mathcal{V}\setminus\mathcal{F}_B
\end{equation}
とする。

hard-fix された $b_j=1$ の寄与を元の syndrome に吸収し、
BP instance $B$ の residual syndrome を
\begin{equation}
    s_a^{(B)}
    =
    s_a
    \oplus
    \bigoplus_{\substack{(j,b_j)\in D_B\\H_{aj}=1}}
    b_j
\end{equation}
と定義する。
したがって、fixed variable は BP message passing から除外され、
その既知の parity contribution のみが residual syndrome に反映される。


\subsection{Step 1: Initial BP}

最初に hard fixation を含まない通常の BP を実行し、
初期 BP instance $B_0$ を得る。

BP が valid correction $\hat{\bm e}$ を生成し、
\begin{equation}
    H\hat{\bm e}=\bm{s}
\end{equation}
を満たす場合には、探索を行う必要はないため直ちに復号を終了する。

BP が収束しない場合には、その posterior LLR の履歴を
以降の探索 guidance として利用する。

この設計の目的は、BP が容易に解ける decoding instance に対して
探索 overhead を発生させないことである。


\subsection{Step 2: Time-Averaged Variable Reliability}

BP instance $B$ における iteration $t$ の posterior LLR を
\begin{equation}
    L_{j,B}^{(t)}
\end{equation}
とする。

loopy BP では LLR が iteration ごとに大きく振動する場合がある。
したがって、最終 iteration の LLR のみを使用せず、
直近 $W$ iterations にわたる clipped average を
\begin{equation}
    \bar L_j^{(B)}
    =
    \frac{1}{W}
    \sum_{t=T-W+1}^{T}
    \operatorname{clip}
    \left(
        L_{j,B}^{(t)},-L_{\mathrm c},L_{\mathrm c}
    \right)
    \label{eq:average_llr}
\end{equation}
と定義する。

ここで $L_{\mathrm c}>0$ は clipping parameter である。
clipping は少数の非常に大きな LLR が reliability score を支配することを防ぐ。

さらに、LLR を $[0,1]$ に正規化した variable confidence を
\begin{equation}
    c_j^{(B)}
    =
    \left|
        \tanh\frac{\bar L_j^{(B)}}{2}
    \right|
    \label{eq:variable_confidence}
\end{equation}
と定義する。
$c_j^{(B)}\simeq 1$ は BP が $x_j$ の値について高い confidence を持つことを、
$c_j^{(B)}\simeq0$ は $x_j$ が曖昧であることを意味する。

Beam Search Decoder~\cite{beamsearch} でも、
複数 iteration にわたる posterior LLR の和の絶対値が
variable reliability として利用されている。
本手法では同じ考え方を採用しつつ、
window averaging、clipping、および $\tanh$ による規格化を行う。


\subsection{Step 3: Check-Level Ambiguity}

variable 単位の ambiguity だけで探索位置を決定すると、
Tanner graph 上の parity constraint の構造を利用できない。
そこで、各 check が現在の BP posterior の下で
どの程度 syndrome constraint と整合しているかを評価する。

LLR convention を
\begin{equation}
    L_j
    =
    \log
    \frac{P(x_j=0)}{P(x_j=1)}
\end{equation}
とすると、
\begin{equation}
    \mathbb{E}[(-1)^{x_j}]
    =
    \tanh\frac{L_j}{2}
\end{equation}
である。

BP posterior の独立近似の下で、
check $a$ が residual syndrome $s_a^{(B)}$ を満たす確率を
\begin{equation}
    Q_a^{(B)}
    =
    \frac{1}{2}
    \left[
        1
        +
        (-1)^{s_a^{(B)}}
        \prod_{j\in\mathcal{N}(a)\cap\mathcal{U}_B}
        \tanh\frac{\bar L_j^{(B)}}{2}
    \right]
    \label{eq:check_probability}
\end{equation}
と定義する。

対応する check ambiguity を
\begin{equation}
    A_a^{(B)}
    =
    1-Q_a^{(B)}
    \label{eq:check_ambiguity}
\end{equation}
とする。

$A_a^{(B)}$ が大きい check は、
現在の BP posterior がその parity constraint と十分整合していない
局所領域を表す。

各 BP instance $B$ について、
\begin{equation}
    \mathcal{C}_B^{(M)}
    =
    \operatorname{TopM}_{a\in\mathcal{C}}
    A_a^{(B)}
\end{equation}
を探索対象 check とする。

ここで重要なのは、
\begin{equation}
    s_a^{(B)}=1
\end{equation}
である residual-unsatisfied check だけでなく、
\begin{equation}
    s_a^{(B)}=0
\end{equation}
である satisfied check も候補に含めることである。

Tesseract~\cite{tesseract} や通常の Decision Tree Decoder
(DTD)~\cite{dtd} では、
unsatisfied check が correction search の起点となる。
本手法ではこれを拡張し、
BP posterior 上で曖昧な satisfied check も
decimation candidate を探索する情報源として利用する。


\subsection{Step 4: Bounded Local Search Around Ambiguous Checks}

\subsubsection{Local variable restriction}

selected check $a\in\mathcal{C}_B^{(M)}$ ごとに、
その近傍の全 variable を探索するのではなく、
confidence の低い $m$ variables のみを候補とする。

\begin{equation}
    \mathcal{V}_{a,B}^{(m)}
    =
    \operatorname{BottomM}_{
        j\in\mathcal{N}(a)\cap\mathcal{U}_B
    }
    c_j^{(B)}.
    \label{eq:local_variables}
\end{equation}

さらに一つの local search における追加 hard fixation の最大数を
$q$ とする。

したがって、$m$ は search width を、
$q$ は search depth を制御する。

一般的な hard-assignment enumeration の最大候補数は
\begin{equation}
    N_{\mathrm{pat}}(m,q)
    =
    \sum_{d=1}^{q}
    {m\choose d}2^d
\end{equation}
であり、$m$ と $q$ によって local search cost を明示的に制限できる。

この段階では別個の max-child や max-leaf parameter は導入せず、
$m$ と $q$ を基本的な探索予算とする。


\subsubsection{Search node}

parent BP instance $B$ から生成された search node $u$ が
新しく追加する hard assignment を
\begin{equation}
    \Delta D_u
\end{equation}
とする。

node $u$ の全 hard assignment は
\begin{equation}
    D_u
    =
    D_B\cup\Delta D_u
\end{equation}
である。

対応する residual syndrome は
\begin{equation}
    s_a^{(u)}
    =
    s_a^{(B)}
    \oplus
    \bigoplus_{\substack{(j,b_j)\in\Delta D_u\\H_{aj}=1}}
    b_j
    \label{eq:node_residual}
\end{equation}
で与えられる。


\subsubsection{Unsatisfied check: solution search}

selected check が
\begin{equation}
    s_a^{(B)}=1
\end{equation}
である場合には、
その check は correction search の起点として利用できる。

この場合には Tesseract~\cite{tesseract} および
DTD~\cite{dtd} と同様に、
partial correction を incrementally に構成する tree search を行う。

最初の branching は selected check $a$ に anchor する。
$\mathcal{V}_{a,B}^{(m)}$ 内の variable を physical weight
$w_j$、次いで index の順に並べ、
canonical branching を構成する。
その後の descendant では residual syndrome に対して
同様の correction-search rule を適用し、
maximum depth $q$ で探索を停止する。

node $u$ までに $1$ に固定された variable の physical cost を
\begin{equation}
    g(u)
    =
    \sum_{\substack{(j,1)\in D_u}} w_j
\end{equation}
とする。

さらに、node $u$ における residual-unsatisfied checks を
\begin{equation}
    \mathcal{R}(u)
    =
    \{a:s_a^{(u)}=1\}
\end{equation}
とする。

free variable $j$ が現在いくつの residual-unsatisfied checks を
同時に cover できるかを
\begin{equation}
    k_j(u)
    =
    \left|
        \{a\in\mathcal{R}(u):H_{aj}=1\}
    \right|
\end{equation}
と定義する。

Tesseract の fractional-cover heuristic に対応して、
\begin{equation}
    h_{\mathrm{frac}}(u)
    =
    \sum_{a\in\mathcal{R}(u)}
    \min_{\substack{
        j\in\mathcal{N}(a)\\
        j\notin\mathcal{F}_u
    }}
    \frac{w_j}{k_j(u)}
    \label{eq:fractional_heuristic}
\end{equation}
を用いる。

solution-search score を
\begin{equation}
    \boxed{
    F_{\mathrm{solve}}(u)
    =
    g(u)+h_{\mathrm{frac}}(u)
    }
    \label{eq:fsolve}
\end{equation}
と定義する。

$g(u)$ は既に採用した fault の physical cost、
$h_{\mathrm{frac}}(u)$ は残りの syndrome を満たすために
最低限必要と推定される追加 cost を意味する。

したがって $F_{\mathrm{solve}}$ は
「この node の descendant に低 physical-cost solution が存在しそうか」
を評価する score である。

Tesseract~\cite{tesseract} では、
この種の $g+h$ score が search ordering 自体を決定する。
本手法でも correction search についてはこの原理を保持する。


\subsubsection{Satisfied check: decimation probe}

selected check が
\begin{equation}
    s_a^{(B)}=0
\end{equation}
の場合、その check を満たすために
必ず新たな fault を追加する必要はない。

したがって satisfied check に対して
Tesseract 型の solution-search tree を構成することには
直接的な論理的根拠がない。

そこで satisfied check は
\emph{decimation probe} としてのみ利用する。

$\mathcal{V}_{a,B}^{(m)}$ 上で depth $q$ 以下の
hard-assignment patterns を生成し、
各 pattern に対して後述する
$F_{\mathrm{guide}}$ のみを計算する。

ここでは、一時的に selected check を unsatisfied にする assignment も
排除しない。
そのような assignment であっても、
他の free variable が後から parity を補償できるためである。


\subsubsection{Guide score}

$F_{\mathrm{guide}}$ の目的は、
node $u$ が valid correction に直接近いかを評価することではない。
その hard assignment を次の BP に与えたときに、
BP の ambiguity を減らす可能性が高いかを評価することである。

まず parent BP posterior の下で、
variable $j$ を $b$ に固定する確率を
\begin{equation}
    P_B(x_j=b)
    =
    \frac{
        1
    }{
        1+\exp[(2b-1)\bar L_j^{(B)}]
    }
\end{equation}
とする。

node $u$ の追加 fixation 数を
\begin{equation}
    d_u=|\Delta D_u|
\end{equation}
とし、
posterior surprisal の平均を
\begin{equation}
    J_{\mathrm{var}}(u|B)
    =
    \frac{1}{d_u}
    \sum_{(j,b)\in\Delta D_u}
    -\log P_B(x_j=b)
    \label{eq:variable_surprisal}
\end{equation}
と定義する。

$J_{\mathrm{var}}$ が小さい pattern は、
parent BP posterior と大きく矛盾しない fixation である。
一方、$\bar L_j\simeq0$ の ambiguous variable では
$b=0,1$ の双方が類似した cost を持つため、
両 branch を自然に探索できる。


次に、hard fixation を適用した場合の ambiguity reduction を
BP を再実行せずに近似する。

parent BP instance の overall ambiguity を
\begin{equation}
    \mathcal{U}(B)
    =
    \frac{1}{|\mathcal{U}_B|}
    \sum_{j\in\mathcal{U}_B}
    \left(1-c_j^{(B)}\right)
    +
    \beta
    \frac{1}{M}
    \sum_{a=1}^{M}
    A_a^{(B)}
    \label{eq:overall_ambiguity}
\end{equation}
とする。

第一項は variable-level ambiguity、
第二項は check-level inconsistency を表す。
$\beta\ge0$ は両者の relative weight である。

candidate node $u$ に含まれる新しい fixed variables は、
仮想的には完全に決定されたものとして confidence $1$ を与える。
free variables の LLR は parent BP の値をそのまま使用する。

また、Eq.~\eqref{eq:node_residual} により residual syndrome を更新し、
free variables のみを使って
\begin{equation}
    \widetilde Q_a(u|B)
    =
    \frac{1}{2}
    \left[
        1
        +
        (-1)^{s_a^{(u)}}
        \prod_{
            j\in\mathcal{N}(a)\setminus\mathcal{F}_u
        }
        \tanh\frac{\bar L_j^{(B)}}{2}
    \right]
\end{equation}
を計算する。

対応する hypothetical ambiguity
$\widetilde{\mathcal U}(u|B)$ は、
Eq.~\eqref{eq:overall_ambiguity} において
新たに fixed された variable の ambiguity を $0$ とし、
$A_a^{(B)}$ を
$1-\widetilde Q_a(u|B)$ に置き換えることで定義する。

decimation による ambiguity reduction を
\begin{equation}
    G_{\mathrm{amb}}(u|B)
    =
    \frac{
        \mathcal U(B)
        -
        \widetilde{\mathcal U}(u|B)
    }{
        d_u
    }
    \label{eq:ambiguity_gain}
\end{equation}
とする。

$d_u$ で割る理由は、
単純に多くの variables を固定した deep node が
常に有利になることを防ぐためである。

最終的に guide score を
\begin{equation}
    \boxed{
    F_{\mathrm{guide}}(u|B)
    =
    J_{\mathrm{var}}(u|B)
    -
    \lambda
    G_{\mathrm{amb}}(u|B)
    }
    \label{eq:fguide}
\end{equation}
と定義する。

ここで $\lambda\ge0$ は guidance strength である。

Eq.~\eqref{eq:fguide} の第一項は
``posterior に照らして不自然な fixation を避ける''
役割を持つ。
第二項は
``その fixation により BP の ambiguity がどの程度解消されそうか''
を評価する。

したがって $F_{\mathrm{guide}}$ が小さいほど、
posterior 上 plausible であり、
かつ BP の曖昧さを効率よく解消すると予測される
decimation pattern である。

なお、hard fixation によって変化する check probability は
固定 variable に隣接する checks のみである。
したがって、実装上は全 $M$ checks の
$\widetilde Q_a$ を再計算する必要はなく、
\begin{equation}
    \mathcal{C}(\Delta D_u)
    =
    \bigcup_{(j,b)\in\Delta D_u}\mathcal{N}(j)
\end{equation}
に対してのみ incremental update を行えばよい。

$F_{\mathrm{guide}}$ は heuristic score であり、
Tesseract の admissible heuristic のような
minimum-weight guarantee を与えるものではない。


\subsection{Step 5: Selection and Execution of Decimated BP Instances}

Step 4 により、一つの parent BP instance から多数の search nodes が得られる。
さらに recursive cycle では複数の parent BP instances が存在する。

cycle $c$ の parent BP beam を
\begin{equation}
    \mathcal{B}_c
    =
    \{B_1,\ldots,B_{K_{\mathrm{keep}}}\}
\end{equation}
とする。

各 $B\in\mathcal{B}_c$ から生成された search nodes を全て統合し、
global candidate pool
\begin{equation}
    \mathcal{P}_c
    =
    \bigcup_{B\in\mathcal{B}_c}
    \mathcal{P}(B)
\end{equation}
を構成する。

ここで、
\begin{equation}
    K_{\mathrm{run}}
\end{equation}
を cycle あたり実際に実行する decimated BP の総数とする。

重要なのは、各 parent BP instance ごとに
$K_{\mathrm{run}}$ 個を実行するのではなく、
全 parent から生成された candidate pool 全体から
合計 $K_{\mathrm{run}}$ 個のみを選ぶことである。
これにより recursive search の計算量の指数的増大を防ぐ。


\subsubsection{Dual-score admission}

correction search と decimation guidance の両方を保持するため、
\begin{equation}
    K_{\mathrm{s}}
    =
    \left\lfloor
    \frac{K_{\mathrm{run}}}{2}
    \right\rfloor,
    \qquad
    K_{\mathrm{g}}
    =
    K_{\mathrm{run}}-K_{\mathrm{s}}
\end{equation}
とする。

$F_{\mathrm{solve}}$ を持つ unsatisfied-check search nodes のうち
上位 $K_{\mathrm{s}}$ 個を
\begin{equation}
    \mathcal{S}_c
    =
    \operatorname{Top}_{K_{\mathrm{s}}}
    \left(
        -F_{\mathrm{solve}}
    \right)
\end{equation}
とする。

また、unsatisfied/satisfied の双方を含む全 candidate から
$F_{\mathrm{guide}}$ 上位 $K_{\mathrm{g}}$ 個を
\begin{equation}
    \mathcal{G}_c
    =
    \operatorname{Top}_{K_{\mathrm{g}}}
    \left(
        -F_{\mathrm{guide}}
    \right)
\end{equation}
とする。

ここでは両 score とも小さい方が良いため、
$\operatorname{Top}(-F)$ と表記している。

初期 selected set を
\begin{equation}
    \mathcal{E}_c
    =
    \operatorname{Unique}
    \left(
        \mathcal{S}_c
        \cup
        \mathcal{G}_c
    \right)
\end{equation}
とする。

両 ranking に同じ decimation pattern が含まれていた場合、
それは一つの BP instance としてのみ実行する。

重複によって
\begin{equation}
    |\mathcal{E}_c|<K_{\mathrm{run}}
\end{equation}
となった場合には、
未選択候補を $F_{\mathrm{guide}}$ の順に追加する。
$F_{\mathrm{guide}}$ 候補を使い切った場合にのみ
$F_{\mathrm{solve}}$ の次順位から補充する。

この guide-priority refill を採用する理由は、
Step 5 の直接的な目的が
``次に実行する BP instance を選ぶこと''
だからである。
$F_{\mathrm{solve}}$ は correction search の観点から重要であるが、
$F_{\mathrm{guide}}$ は BP admission 自体を目的として設計された score である。


\subsubsection{Execution of hard-decimated BP}

選択された各 candidate $u\in\mathcal{E}_c$ に対して、
hard-fixation $D_u$ を課した decimated BP を実行する。

candidate $u$ が parent BP instance $B$ から生成された場合には、
可能な限り $B$ の message state を継承し、
新たに固定された variables とその incident edges のみを
message passing から除外する。

これは毎回 channel prior から BP を cold restart するよりも、
parent BP が既に獲得した局所情報を保持できるためである。

Beam Search Decoder~\cite{beamsearch} でも、
各 path が BP message state を保持し、
branch 後の masked BP に引き継ぐ構造が用いられている。

decimated BP の開始直後および各 iteration 後に
\begin{equation}
    H\hat{\bm e}=\bm{s}
\end{equation}
を確認し、
valid correction が得られた場合には直ちに終了する。


\subsection{Step 6: BP-State Retention and Recursive Cycles}

$K_{\mathrm{run}}$ 個の decimated BP のうち、
valid correction を得られなかった instances について、
次 cycle に残す parent BP instances を選択する。

BP instance $B_u$ の final LLR history から
Eq.~\eqref{eq:average_llr} と
Eq.~\eqref{eq:variable_confidence} を再計算し、
global BP reliability を
\begin{equation}
    \boxed{
    R(B_u)
    =
    \frac{1}{
        |\mathcal{U}_{B_u}|
    }
    \sum_{j\in\mathcal{U}_{B_u}}
    c_j^{(B_u)}
    }
    \label{eq:bp_reliability}
\end{equation}
と定義する。

$R(B_u)\in[0,1]$ であり、
値が大きいほど free variables 全体について
BP posterior が decisive であることを意味する。

Beam Search Decoder~\cite{beamsearch} では、
unmasked variables の accumulated LLR reliability の総和を
path score として用い、
高 score の paths のみを beam width まで保持する。
Eq.~\eqref{eq:bp_reliability} はこの考え方を
normalized confidence として採用したものである。

次 cycle に保持する BP instance 数を
\begin{equation}
    K_{\mathrm{keep}}
    \le
    K_{\mathrm{run}}
\end{equation}
とし、
\begin{equation}
    \boxed{
    \mathcal{B}_{c+1}
    =
    \operatorname{Top}_{K_{\mathrm{keep}}}
    R(B_u)
    }
    \label{eq:beam_retention}
\end{equation}
とする。

したがって一 cycle の情報 flow は
\begin{equation}
\begin{split}
    \{\text{parent BP states}\}
    &\longrightarrow
    \{\text{many cheap search nodes}\}
    \\
    &\longrightarrow
    K_{\mathrm{run}}\text{ selected patterns}
    \\
    &\longrightarrow
    K_{\mathrm{run}}\text{ decimated BP runs}
    \\
    &\longrightarrow
    K_{\mathrm{keep}}\text{ retained BP states}.
\end{split}
\end{equation}

この処理を最大 $C_{\max}$ cycles 繰り返す。

ここで
$K_{\mathrm{run}}$ は一 cycle に評価する BP candidate 数、
$K_{\mathrm{keep}}$ は recursive search の width を表す。

両者を分離することで、
一時的には複数の candidate を評価しつつ、
次 cycle に伝播する state 数を小さく保つことができる。


\subsection{Step 7: OSD Fallback}

最大 $C_{\max}$ cycles を実行しても valid correction が得られない場合には、
OSD fallback を実行する。

最後に保持された BP states のうち
\begin{equation}
    B^\star
    =
    \arg\max_{B\in\mathcal{B}_{C_{\max}}}
    R(B)
\end{equation}
を選択し、
その full signed LLR vector を OSD の reliability input とする。

初期実装では OSD-0 を用いることができる。

OSD が生成した correction は、
それ以前の search/BP solution と同様に
original decoding problem に対して
\begin{equation}
    H\hat{\bm e}=\bm{s}
\end{equation}
を満たすことを最終的に検証する。

retained BP state が存在しない場合には、
physical channel prior を OSD input とする。


\subsection{先行研究との関係}

\paragraph{Tesseract.}
Tesseract~\cite{tesseract} は error-set space 上の
A$^\ast$ search として decoding を定式化し、
physical cost $g$ と residual syndrome に対する
admissible heuristic $h$ を用いる。
本手法の $F_{\mathrm{solve}}$ はこの考え方を継承している。

一方、本手法では探索位置自体を BP posterior から得られる
check ambiguity $A_a$ で決定し、
さらに satisfied checks からも decimation patterns を生成する。
また、search result は correction candidate だけでなく、
次の BP instance を生成する hard-fixation pattern としても利用される。


\paragraph{Decision Tree Decoder.}
DTD~\cite{dtd} では partial correction ごとに cost を与え、
最小 cost の live node を逐次探索する。
Height-bound DTD では syndrome-height lower bound を primary cost とし、
decimated BP の LLR を tie-breaker として使用する。
BP-DTD では decimated BP LLR を直接 tree cost update に使用する。

本手法も BP posterior を search guidance に利用する点では共通するが、
各 search node で BP を再実行しない。
一つの parent BP state から多数の search nodes を
$F_{\mathrm{solve}}$ と $F_{\mathrm{guide}}$ により安価に評価し、
global pool から選択された $K_{\mathrm{run}}$ nodes のみについて
BP を実行する。
したがって、search と BP evaluation を二段階に分離している。


\paragraph{Beam Search Decoder.}
Beam Search Decoder~\cite{beamsearch} は、
BP accumulated LLR から least-reliable variable を選択し、
その variable を $0/1$ に branch した後、
各 branch で masked BP を実行する。
BP 後の path reliability score により beam width 個まで prune する。

本手法は
Eq.~\eqref{eq:average_llr} と
Eq.~\eqref{eq:bp_reliability} において
この temporal-LLR reliability の考え方を採用する。

ただし、branching unit は単一の least-reliable variable ではない。
check ambiguity $A_a$ により局所領域を選択し、
その周囲で bounded correction search と
decimation-pattern search を行う。
さらに、BP 実行前には
$F_{\mathrm{solve}}$ と $F_{\mathrm{guide}}$ による screening を行うため、
生成された全 branch に対して BP を実行する必要がない。


\subsection{アルゴリズム全体の要約}

本手法は次の三段階の score hierarchy を持つ。

\begin{equation}
    \boxed{
    A_a:
    \quad
    \text{どの check 周辺を探索するか}
    }
\end{equation}

\begin{equation}
    \boxed{
    F_{\mathrm{solve}},
    F_{\mathrm{guide}}:
    \quad
    \text{どの decimation pattern に BP を実行するか}
    }
\end{equation}

\begin{equation}
    \boxed{
    R(B):
    \quad
    \text{実行後のどの BP state を次 cycle に残すか}
    }
\end{equation}

したがって全体の処理は
\begin{equation}
\boxed{
\begin{aligned}
\text{BP}
&\rightarrow
\text{time-averaged LLR}
\\
&\rightarrow
\text{check ambiguity selection}
\\
&\rightarrow
\text{bounded local search}
\\
&\rightarrow
\text{dual-score candidate selection}
\\
&\rightarrow
K_{\mathrm{run}}\text{ decimated BPs}
\\
&\rightarrow
K_{\mathrm{keep}}\text{ BP-state beam}
\\
&\rightarrow
\text{recursive search}
\\
&\rightarrow
\text{OSD fallback}.
\end{aligned}
}
\end{equation}

この構造では、correction search と BP decimation search を
同一の search nodes 上で実行しつつ、
両者の目的関数を意図的に分離している。
$F_{\mathrm{solve}}$ は低 physical-cost solution の発見を目的とし、
$F_{\mathrm{guide}}$ は BP ambiguity の解消を目的とする。
その後、実際に BP を実行した結果は
$R(B)$ によって独立に評価される。

この三段階の分離により、
search score の誤差がそのまま recursive BP beam の選択誤差になることを避け、
``solution を探す search'' と
``次に有効な decimation を探す search''
を一つの decoder 内で共存させる。


\begin{thebibliography}{9}

\bibitem{tesseract}
L.~A.~Beni, O.~Higgott, and N.~Shutty,
``Tesseract: A Search-Based Decoder for Quantum Error Correction,''
arXiv:2503.10988, 2025.

\bibitem{dtd}
K.~R.~Ott, B.~Het\'enyi, and M.~E.~Beverland,
``Decision-Tree Decoders for General Quantum LDPC Codes,''
arXiv:2502.16408, 2025.

\bibitem{beamsearch}
M.~Ye, D.~Wecker, and N.~Delfosse,
``Beam Search Decoder for Quantum LDPC Codes,''
arXiv:2512.07057, 2025.

\end{thebibliography}